\chapter{Công trình liên quan}

Chương này tập trung phân tích các công trình nghiên cứu khoa học và các hệ thống phần mềm đã có liên quan đến lĩnh vực thư viện số và hệ thống gợi ý. Từ đó, chương sẽ chỉ ra những khoảng trống nghiên cứu và định vị hướng đi của đề tài.

\section{Các công trình nghiên cứu khoa học}

Qua việc tổng hợp các công trình nghiên cứu, có thể thấy các hướng tiếp cận phổ biến trong việc xây dựng hệ thống gợi ý sách cho thư viện có những đặc điểm sau:

\begin{itemize}
    \item \textbf{Ưu điểm và xu hướng chung:}
    \begin{itemize}
        \item Hầu hết các mô hình gợi ý đều theo hướng \textbf{hybrid}, kết hợp giữa Lọc dựa trên nội dung (Content-Based) và Lọc cộng tác (Collaborative Filtering) để cân bằng giữa độ cá nhân hóa và khả năng khai thác hành vi người dùng.
        \item Một số nghiên cứu cải tiến kết quả bằng cách thêm các cơ chế xếp hạng (ranking) hoặc gợi ý theo cụm người dùng (clustering), đặc biệt hiệu quả trong môi trường dữ liệu vừa và nhỏ.
        \item Việc đánh giá mô hình thường dựa trên dữ liệu thực tế từ các thư viện, giúp kiểm chứng rõ ràng hiệu quả của thuật toán.
    \end{itemize}

    \item \textbf{Hạn chế phổ biến:}
    \begin{itemize}
        \item Nhiều công trình tập trung sâu vào phần mô hình thuật toán nhưng ít chú trọng đến việc triển khai thành một hệ thống hoàn chỉnh với đầy đủ backend và frontend.
        \item Các hướng đi hiện đại như tìm kiếm ngữ nghĩa (semantic search) hay các mô hình Transformer tuy tiềm năng nhưng đòi hỏi chi phí tính toán và nguồn lực lớn, vượt quá khả năng của nhiều dự án quy mô nhỏ.
    \end{itemize}
\end{itemize}

\section{Các hệ thống và phần mềm tương tự}

Hiện nay, có nhiều hệ thống quản lý thư viện điện tử đã được triển khai rộng rãi, tiêu biểu có thể kể đến:

\begin{itemize}
    \item \textbf{Opac-HCMUT-VNU:} Là hệ thống của Đại học Quốc Gia TP.HCM, đáp ứng tốt các nghiệp vụ tra cứu và mượn-trả cơ bản. Tuy nhiên, chức năng tìm kiếm vẫn chủ yếu dựa vào từ khóa và lọc và chưa tích hợp tính năng gợi ý thông minh, tìm kiếm theo ngữ nghĩa.

    \item \textbf{Koha:} Là một phần mềm quản lý thư viện mã nguồn mở rất phổ biến, cung cấp đầy đủ chức năng quản lý và báo cáo. Mặc dù vậy, giao diện của Koha đã cũ, chưa hỗ trợ gợi ý tự động và chưa có cơ chế phân tích hành vi người dùng.

    \item \textbf{Aleph (Ex Libris):} Đây là một hệ thống thương mại quy mô lớn, tích hợp nhiều cơ sở dữ liệu học thuật. Nhược điểm chính của hệ thống này là chi phí triển khai cao và không phù hợp cho các thư viện quy mô nhỏ như cấp khoa hoặc trường đại học nhỏ.
\end{itemize}

\section{Kết chương và xác định khoảng trống nghiên cứu}

Từ việc phân tích các công trình khoa học và hệ thống thực tế, nhóm nhận thấy một số \textbf{khoảng trống (gaps)} quan trọng mà đề tài này hướng đến giải quyết:

\begin{enumerate}
    \item \textbf{Khoảng trống giữa Mô hình và Thực thi:} Nhiều nghiên cứu tập trung vào việc xây dựng các mô hình gợi ý phức tạp nhưng lại bỏ qua việc triển khai chúng thành một hệ thống phần mềm hoàn chỉnh, có tính ứng dụng thực tế. Đa số các mô hình hiện nay tập trung nhiều vào thuật toán gợi ý, trong khi ít chú trọng đến việc triển khai thành hệ thống hoàn chỉnh.

    \item \textbf{Khoảng trống về Trải nghiệm Người dùng:} Các hệ thống hiện tại, kể cả mã nguồn mở như Koha, thường có giao diện cũ, chưa tối ưu cho các thiết bị hiện đại và thiếu các tính năng thông minh giúp nâng cao trải nghiệm tìm kiếm, khám phá tài liệu của người dùng.

    \item \textbf{Khoảng trống về Tính khả thi:} Các giải pháp gợi ý tiên tiến thường yêu cầu nguồn dữ liệu lớn hoặc xử lý phức tạp, vượt quá phạm vi và nguồn lực của một nhóm sinh viên.
\end{enumerate}

\textbf{Hướng đi của dự án:} Đề tài này không đặt nặng việc phát triển một mô hình AI quá chuyên sâu. Thay vào đó, dự án tập trung vào việc \textbf{bắc cầu cho các khoảng trống trên} bằng cách:
\begin{itemize}
    \item Xây dựng một \textbf{hệ thống hoàn chỉnh} với kiến trúc frontend-backend tách biệt, dễ bảo trì và mở rộng.
    \item Ưu tiên \textbf{trải nghiệm người dùng} với giao diện hiện đại và tích hợp các tính năng gợi ý, tìm kiếm thông minh một cách liền mạch.
    \item Áp dụng một \textbf{giải pháp AI khả thi} (mô hình hybrid cơ bản) để tạo ra một sản phẩm hoạt động ổn định, có giá trị thực tiễn và phù hợp với bối cảnh của một đồ án tốt nghiệp.
\end{itemize}
